\chapter{Methodology}
\label{ch:methodology}

% --------------------------------------------------------------------------
\section{Research Questions}
\label{sec:rq}
% --------------------------------------------------------------------------

This thesis is structured around three research questions.

\textbf{RQ1.} Can a 1.5-billion-parameter code language model learn to generate
syntactically and semantically valid eBPF programs through supervised fine-tuning
on a labeled dataset of program--verifier-log pairs collected from the BPF
verifier?

\textbf{RQ2.} Can a GRPO reinforcement-learning loop with live KCOV kernel PC
coverage as the reward signal increase verifier path coverage beyond what the SFT
model achieves at inference time?

\textbf{RQ3.} What failure modes arise in GRPO training with live kernel feedback,
and how can they be diagnosed and systematically addressed?

RQ1 is answered by the SFT experiment (Chapter~\ref{ch:results},
Section~\ref{sec:results-sft}).
RQ2 is partially answered by RL run~1 and is the subject of the planned run~2.
RQ3 is the central analytical contribution of this thesis: the plateau diagnosis,
the root-cause identification, and the reward redesign all address this question.

% --------------------------------------------------------------------------
\section{Data Collection: Buzzer and the KCOV Discovery}
\label{sec:buzzer-data}
% --------------------------------------------------------------------------

The SFT dataset requires a large corpus of BPF programs paired with their
verifier-log outputs.
No such public corpus existed; this work collects one by instrumenting buzzer at
the kernel FFI boundary.

Buzzer is a Go binary that generates eBPF programs using one of two built-in
strategies and submits each program to the kernel via a function called
\texttt{ValidateEbpfProgram}, which wraps the \texttt{BPF\_PROG\_LOAD} system
call and returns the verifier verdict and log.
The \texttt{coverage\_based} strategy also opens the KCOV device, arms tracing
before each \texttt{BPF\_PROG\_LOAD}, and collects the resulting PC counts,
which it exposes as human-readable metrics via an internal HTTP server on a
per-runner port.

A close reading of the buzzer source code revealed that these KCOV metrics are
display-only.
The \texttt{coverage\_based} strategy generates programs using the same
hard-coded Go construction logic regardless of which kernel paths previous
programs have covered.
KCOV data is never fed back to the generation loop.
Buzzer is therefore not a coverage-guided fuzzer in the AFL sense:
the feedback loop that would make it one is absent.
This finding --- that a gap between data collection and mutation existed and was
not previously documented --- was the concrete motivation for building the
closed loop that this thesis implements.

To produce the training corpus, the file
\texttt{fuzzing/buzzer/pkg/units/ffi.go} was modified (approximately 50 lines)
to intercept every call to \texttt{ValidateEbpfProgram} and write a JSONL record
containing the bytecode in hexadecimal, the full verifier log, the
accept/reject verdict, and (where present) the error line and error reason.
Records are appended to a file on a virtio-9p share visible from the host.

The collection infrastructure consists of a single QEMU virtual machine running
a Linux kernel compiled with KASAN only (\texttt{bzImage\_kasan}), managed
by a Docker container that auto-restarts the VM on crash.
Collection ran for approximately one week and produced roughly 2~million
(bytecode, verifier-log) pairs.
Thirty kernel-panic and KASAN-report dumps were captured in
\texttt{results/crash\_logs/} during the collection period.

% --------------------------------------------------------------------------
\section{SFT Pipeline Design}
\label{sec:sft-design}
% --------------------------------------------------------------------------

\subsection{Dataset Curation}
\label{sec:dataset}

The 2-million-entry raw corpus has a highly skewed error-class distribution:
a small number of verifier error types (such as ``invalid access to map value'')
account for the majority of rejected programs, while dozens of rarer error classes
are represented by only a handful of examples each.
Training on the raw distribution would bias the model toward generating programs
that trigger the most common errors and would provide little signal for the
rarer, potentially more interesting error conditions.

The curation procedure applies three steps.
First, verifier error text is normalised: numeric addresses, instruction offsets,
and register values are masked to canonical tokens so that distinct error messages
that differ only in runtime-specific numbers are mapped to the same class label.
This grouping exposes the underlying error taxonomy, yielding 26 distinct error
classes.
Second, a per-class cap of 2\,000 samples is applied to the rejected-program
entries, preventing any single error class from dominating the training
distribution.
Third, all 13\,153 valid programs are retained without capping: because valid
programs are already a minority in the raw corpus, including them in full is
necessary to give the model sufficient exposure to the accepted-program
distribution that the SFT objective asks it to learn.

The final dataset contains 27\,514 examples: 13\,153 valid programs and 14\,361
invalid programs spread across the 26 error classes.
The dataset is split 90/10 for training and evaluation.
It is published on HuggingFace as \texttt{Strhata/ebpf-corpus} and serves as the
SFT training corpus for this thesis.

\subsection{Representation: Assembly Format vs.\ Raw Bytecode}
\label{sec:format-decision}

BPF programs can be represented at two levels: as raw binary bytecode
(8-byte little-endian structs, one per instruction) or as verifier-log assembly
(the textual disassembly that the kernel's verifier log prints).
The choice of representation for the LLM's output vocabulary is consequential for
training quality and is justified here from first principles.

Raw bytecode is an opaque byte sequence with no natural-language grounding.
Each byte position within an instruction struct carries a different semantic role
--- opcode, destination register packed with source register, 16-bit signed
offset, 32-bit immediate --- but this role is invisible from the token itself.
A language model tokenizing raw bytes would need to learn these positional
semantics entirely from the statistical structure of the training data, with no
character- or word-level surface cues to anchor the learning.

Verifier-log assembly makes these semantics explicit.
Each instruction is a human-readable line of the form \texttt{N: (XX) text},
where \texttt{N} is the instruction index, \texttt{XX} is the two-hex-digit
opcode byte, and \texttt{text} is a natural-language-style mnemonic with named
register operands.
Register names (r0--r9 for general-purpose registers, r10 for the frame pointer)
match the identifiers that appear in the accompanying verifier-log error messages,
making the connection between instruction operands and verifier error text
transparent at the token level.
Jump targets are expressed as relative offsets labelled by the verifier
(\texttt{+N} or \texttt{-N}), which carry instruction-level semantic content that
an absolute bytecode offset does not.

Two further practical advantages reinforce the choice.
The verifier-log format is the exact output of buzzer's data extraction path,
requiring no intermediate conversion at training time.
And because the format includes the opcode byte \texttt{(XX)} explicitly, the
pure-Python encoder described in Chapter~\ref{ch:implementation} can reconstruct
the raw bytecode from model output deterministically, without depending on an
assembler or a clang installation.

An example of the format as it appears in model output is:

\begin{lstlisting}[caption={Verifier-log assembly format --- five instructions}]
0: (b7) r1 = 0
1: (63) *(u32 *)(r10 -4) = r1
2: (bf) r2 = r10
3: (07) r2 += -4
4: (95) exit
\end{lstlisting}

% --------------------------------------------------------------------------
\section{Reinforcement Learning Design}
\label{sec:rl-design}
% --------------------------------------------------------------------------

\subsection{Reward Function v1}
\label{sec:reward-v1}

RL run~1 uses GRPO with $G = 4$ completions per prompt, KL penalty
$\beta = 0.01$, and a maximum completion length of 600 tokens.
These hyperparameters reflect the memory budget of the local training GPU
(RTX~4070 Laptop, 8.59~GB VRAM): larger $G$ or longer sequences exhaust memory
before a full forward pass.

The reward function assigns a scalar score to each generated program after it has
been submitted to the \texttt{kcov\_validator} binary running in the KCOV-instrumented
VM.
Five tiers are defined, in order of priority:

\begin{table}[h]
\centering
\caption{Reward tiers for RL run~1}
\label{tab:reward-v1}
\begin{tabular}{llr}
\toprule
Tier & Condition & Reward \\
\midrule
Crash or timeout   & SSH timeout; VM may have crashed                & 2.0 \\
New kernel PCs     & ACCEPTED and at least one PC not in frontier    & 1.0 \\
Valid, no new PCs  & ACCEPTED, all PCs already in frontier           & 0.1 \\
Verifier rejected  & REJECTED by the verifier (any reason)           & 0.1 \\
Encode failure     & Output is not parseable BPF assembly            & 0.0 \\
\bottomrule
\end{tabular}
\end{table}

The crash/timeout tier is assigned the highest reward because an SSH timeout
during program load suggests the VM crashed, which is the strongest evidence of
a verifier defect.
New kernel PCs are the primary exploration signal: reaching a code path not seen
in any previous batch is rewarded at 1.0.
The ``valid, no new PCs'' and ``verifier rejected'' tiers share the same value
(0.1) deliberately: rejection is informative about the program's syntax, not
about exploration quality, and penalising rejection as a hard negative signal
would push the model away from programs that stress verifier error paths ---
exactly the programs of interest.
Encode failures receive zero reward because the model output is unparseable and
provides no coverage information.

The cumulative coverage frontier is maintained as a Python set of kernel PC
addresses, updated after each batch of $G \times B$ programs (where $B$ is the
number of prompts per batch).
A snapshot of the frontier is taken before each batch is evaluated so that all
$G$ completions for a given prompt compete against the same baseline, independent
of evaluation order within the batch.

\subsection{Root Cause: KCOV Bug and \texttt{reward\_std} Collapse}
\label{sec:rl-plateau}

Run~1 (\texttt{rl\_grpo\_v2}) executed 8,370 training steps.
Cumulative new-PC growth ceased at approximately step~1,300, after which the
model made no further progress despite 7,000 additional steps of compute.
Post-run analysis of \texttt{results/grpo\_completions.log} revealed that
\texttt{reward\_std}, computed per batch as the standard deviation of the $G$
rewards, dropped to exactly zero from step~$\approx$1,300 onward.

A zero \texttt{reward\_std} is a definitive indicator of training collapse.
The GRPO advantage for completion~$i$ is
$A_i = (r_i - \bar{r}) / \sigma_r$ (Equation~\ref{eq:grpo-advantage}).
When $\sigma_r = 0$ --- which occurs when all $G$ completions in a group receive
identical rewards --- every advantage is zero regardless of the reward values,
the policy gradient is identically zero, and no parameter update occurs.
The model is frozen in place while the training loop continues consuming compute
and GPU time.
This condition, known as gradient starvation, is analysed formally
in~\cite{grpo-starvation}.

The root cause was a conditional in the \texttt{kcov\_validator} Go binary:
\texttt{readPCs()} was called only inside a branch taken when
\texttt{verdict == ACCEPTED}.
For every REJECTED program, the function returned \texttt{pcs:~[]}.
At run~1 inference time, the SFT model was still generating a high proportion of
rejected programs.
Consequently, the majority of completions in each GRPO group received reward~0.1
with an empty PC set; every completion appeared identical from the reward
function's perspective; and within-group variance collapsed to zero.

A secondary hypothesis, which cannot be verified from run~1 data, concerns the
prompt template.
The SFT training data used a \texttt{Status:~VALID} header as the prompt, and
the same prompt was reused verbatim during RL inference.
This may have anchored the model's generation distribution toward valid programs,
reducing the diversity of the $G$ completions within each group and thereby
contributing to reward homogeneity.
Because the KCOV bug and the prompt-bias confound coexisted in run~1, they cannot
be disentangled without a second run using a neutral prompt.

The bug was fixed in commit \texttt{c8a9d02}: the \texttt{readPCs()} call was
moved outside the \texttt{if verdict == ACCEPTED} branch, so both ACCEPTED and
REJECTED programs now return their full KCOV trace.

\subsection{Reward Redesign: Depth-Based Verdict-Blind}
\label{sec:reward-redesign}

The diagnosis in Section~\ref{sec:rl-plateau} identifies the structural
requirement for GRPO training to proceed: the reward function must produce
non-zero within-group variance $\sigma_r^2 > 0$ even when all $G$ completions
are rejected by the verifier.
The v1 reward fails this requirement because REJECTED and accepted-but-no-new-PCs
programs both receive reward~0.1, and when the model generates mostly rejected
programs these values are all equal.

The redesigned reward removes the accept/reject verdict from the reward formula
entirely.
A program is scored by two components: a depth score proportional to the number
of kernel PCs reached during verification, and a discovery bonus for PCs not seen
in any prior batch.
Both components are defined exclusively in terms of KCOV traces, which are now
available for rejected programs following the bug fix.

Formally, let $|\mathrm{pcs}|$ be the number of distinct kernel PCs returned by
\texttt{kcov\_validator} for a given program, $M$ the largest PC count observed
across all programs evaluated so far (updated after each batch), and
$S_{\mathrm{snap}}$ a snapshot of the cumulative coverage frontier taken before
the current batch is evaluated.
The reward is:
\begin{align}
  \mathrm{depth} &= \min\!\left(0.5,\;
                     \frac{|\mathrm{pcs}|}{M} \cdot 0.5\right), \label{eq:depth}\\
  \mathrm{bonus} &= 1.0 \text{ if } \exists\, p \in \mathrm{pcs} : p \notin S_{\mathrm{snap}},
                   \text{ else } 0, \label{eq:bonus}\\
  r              &= \mathrm{depth} + \mathrm{bonus}. \label{eq:reward-v2}
\end{align}

The pre-batch snapshot $S_{\mathrm{snap}}$ is critical for fairness: it ensures
that all $G$ completions for a given prompt compete against the same coverage
frontier regardless of the order in which they are evaluated within the batch.
The depth component is capped at 0.5 to prevent a single program with an
unusually large PC trace from saturating the reward scale and dominating the
group-relative advantage.
The discovery bonus is capped at 1.0; combined with the depth cap, the maximum
reward per program is 1.5 under normal operation.

Two special cases override the formula: an encode failure returns 0.0 (the
program is not parseable and provides no coverage information), and an SSH
timeout or VM crash returns 2.0 (the highest reward, preserving the crash signal
from v1 as the most valuable exploration outcome).

The key property of this design is that within-group reward variance is
structurally guaranteed even in a fully rejected group.
Programs that trigger deeper verifier paths --- more complex type-checking,
longer control-flow analysis --- produce longer KCOV traces and therefore receive
higher depth scores, providing a gradient signal proportional to exploration
depth regardless of whether the verifier ultimately accepts or rejects the
program.

The redesigned reward function is implemented in \texttt{ml/reward.py}.
A full training run (run~2) with this reward and a neutral prompt is left as
future work and is described in Section~\ref{sec:run2}.
